\clearpage
\scp{\tsc{Southwest Tanimbar}}{04-07}
\tsc{Southwest Tanimbar} contains Makatian
and Seluwasan, spoken on the southwest part of the island of Tanimbar,
as well as Selaru spoken on an island of the same name immediately
to the southwest of Tanimbar.
These languages are placed within \tsc{Timor-Babar} as they have
*p > {\0} word initially and medially.
Examples in addition to those in
\trf{04-01} (\prf{tab:04-01}) are given in \trf{04-62}.

\begin{table}\tcp{\tsc{Southwest Tanimbar} *p > {\0} /{\#}{\gap}, /V{\gap}V}{04-62}
\begin{adjustbox}{width=\textwidth}
	%\begin{threeparttable}
	\st{2pt}\begin{tabular}{lllllllll}\lsptoprule
local gl.	&	stingray		&	turtle			&	choose			&	belly button	&	arm span		&	hungry			&	dream					&	calcium		\\	
PMP				&	*\tbr{p}aRih&	*\tbr{p}əñu	&	*\tbr{p}iliq&	*\tbr{p}usəj	&	*də\tbr{p}a	&	*la\tbr{p}aR&	*ma-hi\tbr{p}i&	*qa\tbr{p}uR	\\	\midrule
Selaru		&	{\hr}ary		&	{\hr}enw		&	\hp{n}-ily	&	{\hr}usa			&	{\hr}re			&	\hp{na}-mlar&	-mey					&	{\hr}aur		\\	
Seluwasan	&	{\hr}ay-e 	&	{\hr}enw-e 	&	\hp{n}-idy	&	{\hr}uska-ne	&	{\hr}rey-e	&	na-mlay 		&	{\hr}moy 			&	{\hr}aw-e 	\\	
Makatian	&	{\hr}ay-e		&	{\hr}enw-e	&	n-iʤ				&	{\hr}uska-ne	&	{\hr}rey-e	&	na-mlay			&	{\hr}moy 			&	{\hr}aw-e 	\\	
		\lspbottomrule
	\end{tabular}
		%\begin{tablenotes}\tnsize
			%\item[†]
			%%\item[‡]
			%%\item[Δ]
			%%\item[‖]
			%%\item[¶]
		%\end{tablenotes}
	%\end{threeparttable}
\end{adjustbox}
\end{table}

%Yamdena fari, feni, -fily, fusr-e, def-e, lafar, mefi, yafur, (kekap), ?


While the \tsc{Southwest Tanimbar} languages have initial
and medial *p > {\0} initially,
word finally *p = \it{p} in Makatian
and Seluwasan, and *p > \it{f,
h} /{\gap}{\#} in Selaru.
Examples are given in \trf{04-63}.\footnote{
		Selaru /f/ has multiple allophones: /f/ → [f] {\tl} [h] /{\#}{\gap}
		and /f/ → [f] {\tl} [ɸ] {\tl} [h] {\gap}{\#}.
		Labio-dental [f] is most common word finally \citep[15]{CowardCoward2000}.
		Despite the realisation of Selaru /f/ as [h],
		/h/ is also an independent phoneme in Selaru.
		\citet[15]{CowardCoward2000} note “[{\ldots}] there may be social register/dialectal
		variations associated with this phoneme.
		In the more isolated Selaru villages,
		[f] is used more frequently than [h].”}

The reflexes of final *p in Makatian
and Seluwasan provide evidence that we must also reconstruct
**p word finally in at least some words at the level of Proto-Timor-Babar.
Additionally, as discussed at the beginning of this chapter,
it is likely that \tsc{Timor-Babar} *p > **h went through an
intermediate labial fricative stage **f or **ɸ.
This intermediate stage is retained word finally in Selaru.

\begin{table}\tcp{\tsc{Southwest Tanimbar} *p = **p /{\gap}{\#}}{04-63}
%	\begin{threeparttable}\st{6pt}
	\begin{tabular}{lllll}\lsptoprule
local gl.	&	laugh	&	fly (v.)	&	live	&	shine\\
PMP	&	**mali\tbr{p}	&	*laya\tbr{p}	&	*ma-qudi\tbr{p}	&	*dila\tbr{p}\\ \midrule
Selaru	&	\hp{n-}ma\tbr{h}is (\it{met.})	&	\cg{(-nem)}	&	\hp{*n}-mori\tbr{f}	&	-dela\tbr{f}\\
Seluwasan	&	\hp{n-}madu\tbr{p}	&	{\hr}sdyu\tbr{p}	&	{\hr}n-moru\tbr{p}	&	\cg{(elan)}\\
Makatian	&	n-maʤu\tbr{p}	&	{\hr}n-ʤi\tbr{p}	&	{\hr}n-moru\tbr{p}	&	\cg{(n-tal)}\\
		\lspbottomrule
	\end{tabular}
		%\begin{tablenotes}\tnsize
			%\item[†] {Selaru \it{-mahis} ‘laugh’ has consonant metathesis. from earlier **masih.}
		%\end{tablenotes}
	%\end{threeparttable}
\end{table}

Evidence against placing these languages in \tsc{Seram-Tanimbar-Bomberai},
to which their neighbours Yamdena
and Fordata belong, comes from not sharing the merger of *z/*d,
which is one of the two main defining sound changes of
\tsc{Seram-Tanimbar-Bomberai} (see \crf{ch:STB}).
Instead, \tsc{Southwest Tanimbar} languages have *z > \it{sy, s}
and *d > \it{r} in most cases,\footnote{
		There are two cases of *d = \it{d}
		in the available data: *\tbr{d}əŋəR > Makatian, Seluwasan \it{na-\tbr{d}ena} ‘hear’
		(alongside \it{n-ren} `hear')
		and *\tbr{d}iri > Selaru \it{-mdiry} `stand'.
		Word finally *d > \it{t} in Makatian
		and Seluwasan; *lahu\tbr{d} > \it{lo\tbr{t}-e} ‘sea’, *sumaŋə\tbr{d} >
		Makatian \it{smwan\tbr{t}-ane}, Seluwasan \it{smwan\tbr{t}-ane} `spirit',
		and *tuhu\tbr{d} > Seluwasan \it{tu\tbr{t}-e} `knee'.
		The only reflex of word-final *d in Selaru identified
		is *sumaŋəd > \it{smwakna-} which appears to have *d > **t > \it{k}.}
as shown in \trf{04-64} and \trf{04-65} respectively.

\begin{table}\tcp{\tsc{Southwest Tanimbar} *z > \it{s}}{04-64}
	\begin{threeparttable}\st{6pt}
	\begin{tabular}{llllll}\lsptoprule
local gl.	&	far										&	path							&	rain							&	ladder								&	green	\\
PMP				&	*\tbr{z}auq						&	*\tbr{z}alan			&	*qu\tbr{z}an			&	*haRə\tbr{z}an\su{†}	& **mo\tbr{z}o\su{‡}	\\ \midrule
Selaru		&	{\hr}\tbr{s}o{\tl}so	&	{\hr}\tbr{s}al		&	{\hr}u\tbr{s}			&	{\hr}are\tbr{s}				& {\hr\hr}mo\tbr{l}{\tl}mol\\
Seluwasan	&	{\hr}\tbr{sy}oan			&	{\hr}\tbr{sy}alan	&	{\hr}wi\tbr{s}in	&	{\hr}e\tbr{s}an				& \\
Makatian	&	{\hr}\tbr{sy}oen			&	{\hr}\tbr{sy}alan	&	{\hr}wi\tbr{s}in	&	{\hr}e\tbr{s}an				& {\hr\hr}mo\tbr{s}{\tl}mos-e\\
	\lspbottomrule
	\end{tabular}
		\begin{tablenotes}\tnsize
			\item[†] {Makatian and Seluwasan preserve the second
								two syllables of *haRəzan, while Selaru preserves the first two.}
			\item[‡] {Cognates of **mozo ‘green’ occur widely in other \tsc{Timor-Babar} languages.
								See the note to \trf{04-27} (\prf{tab:04-27}) and \srf{15-04-02-01} for additional examples.
								Selaru \it{molmol} `green, blue' has irregular *z > \it{l},
								which may point to this word being a loan, but no forms with
								a medial liquid are located nearby.}
		\end{tablenotes}
	\end{threeparttable}
\end{table}

\begin{table}\tcp{\tsc{Southwest Tanimbar} *d > \it{r}}{04-65}
%\begin{adjustbox}{width=\textwidth}
	\begin{threeparttable}\st{3pt}
	\begin{tabular}{llllllll}\lsptoprule
local gl.	&	inside	&	blood	&	bone	&	dugong	&	filth	&	long	&	behind\\
PMP	&	*\tbr{d}aləm	&	*\tbr{d}aRaq	&	*\tbr{d}uRi	&	*\tbr{d}uyuŋ	&	*\tbr{d}aki	&	*ana\tbr{d}uq	&	*ma-u\tbr{d}əhi\\ \midrule
Selaru	&	{\hr}k\tbr{r}ala	&	{\hr}\tbr{l}ara\su{‡}	&	{\hr}\tbr{l}ury\su{‡}	&	{\hr}\tbr{r}ui	&		&	\hp{*a}na\tbr{r}{\tl}narw	&	{\hr}\cg{(lyaw)}\\
Seluwasan	&	{\hr}la\tbr{r}um\su{†}	&	{\hr}\tbr{r}aw-e	&	{\hr}\tbr{r}uy-e	&	{\hr}\tbr{r}uin-e	&	{\hr}\tbr{r}ay-e	&	\hp{*a}na\tbr{r}un	&	{\hr}lmu\tbr{r}\\
Makatian	&	{\hr}\tbr{r}yalam	&	{\hr}\tbr{r}aw-e	&	{\hr}\tbr{r}uy-e	&	{\hr}\tbr{r}un-e	&	{\hr}\tbr{r}ay-e	&	\hp{*a}na\tbr{r}un	&	{\hr}lmu\tbr{r}\\
		\lspbottomrule
	\end{tabular}
		\begin{tablenotes}\tnsize
			\item[†] {Makatian \it{larum} inside has consonant metathesis from **ralum.}
			\item[‡] {Selaru \it{lara} `blood' and \it{lury} `bone' have regressive
								dissimilation of initial **r.
								This also occurs in nearby \tsc{Nuclear Tanimbar-Bomberai}
								languages (\srf{11-05-01}), but not neighbouring Yamdena.}
		\end{tablenotes}
	\end{threeparttable}
%\end{adjustbox}
\end{table}

The strongest evidence for \tsc{Southwest Tanimbar} as a distinct
node comes from the changes affecting PMP *l.
Both \citet[249]{Mills1991} and \citet[210f]{Blust2006} have
pointed out that Selaru has a split of *l > \it{s,
l} with the split conditioned by the quality of adjacent vowels.

\newpage
The Makatian and Seluwasan data provides more context for understanding this
change, and allows us to identify the changes affecting *l as an
exclusively shared innovation at the level of Proto-Southwest Tanimbar.
The reflexes of *l in \tsc{Southwest Tanimbar} are summarised in \trf{04-66}.

\begin{table}
	\centering\tcp{Reflexes of *l in \tsc{Southwest Tanimbar}}{04-66}
	\wg\st{6pt}\begin{tabular}{lllll}\lsptoprule
PMP	&	Makatian	&	Seluwasan	&	Selaru	&	env.	\\		\midrule
*l	&	ʤ					&	d					&	s				&	/{\gap}V\tsc{[+high]}	\\
*l	&	ʤ					&	ʤ					&	s				&	/{\gap}e	\\	
*l	&	ʤ					&	ʤ					&	l				&	/V\tsc{[+high]}{\gap}	\\	
*l	&	ʤ					&	d					&	l				&	/C{\gap}V\tsc{[+high]}	\\	
*l	&	ʤ					&	d					&	l, s		&	/{\gap}C\tsc{[+glide]}	\\	
*l	&	l					&	l					&	l				&	/{\gap}a,o\\	
		\lspbottomrule
	\end{tabular}
\end{table}

Before high vowels and \it{e} (from *ə) *l > \it{ʤ}
in Makatian and *l > \it{s} in Selaru.
In Seluwasan *l > \it{ʤ} before \it{e}
and *l > \it{d} before high vowels.
Examples word initially are given in \trf{04-67}
and word medially in \trf{04-68}.\footnote{
		*quliŋ `rudder' > Makatian, Seluwasan \it{wiʤ-e},
		Selaru \it{wilin} has irregular
		*l = \it{l}.
		*u > \it{wi} is also irregular in Selaru,
		pointing to this word likely being a secondary loan.
		Teun-Nila-Serua also have similar irregular
		reflexes of this word (see \trf{04-39}, \prf{tab:04-39}).
		}

\begin{table}\centering\tcp{\tsc{Southwest Tanimbar} *l > \it{s -- d, ʤ -- ʤ} /{\#}{\gap}\it{u, i, e}}{04-67}
\begin{adjustbox}{width=\textwidth}
	\begin{threeparttable}\st{2pt}
	\begin{tabular}{lllllllll}\lsptoprule
local gl.	&	five					&	fly (v.)		&	centipede				&	tear (n.)					&	taboo							&	sun, day				&	ginger					&	mortar	\\	
PMP				&	*\tbr{l}ima		&	*\tbr{l}ayap&	*qa\tbr{l}u-hipan\su{†}&*\tbr{l}uhəq&	**\tbr{l}uli\su{‡}&	*qa\tbr{l}əjaw	&	*\tbr{l}aqia		&	*\tbr{l}əsuŋ	\\	\midrule
Selaru		&	{\hr}\tbr{s}im&	\cg{(-nem)}	&	\tcg{(lolohisi)}&	{\hr}\tbr{s}ua		&	(h)a\tbr{s}usy		&	\tbr{s}ew\su{Δ}	&									&	\cg{(nesw)}\su{‖}	\\	
Seluwasan	&	{\hr}\tbr{d}im&	s\tbr{d}yup	&	\tbr{d}ian-e		&	{\hr}\tbr{d}iwa-n	&	\tbr{d}udy				&	\tbr{ʤ}ew-e			&	{\hr}\tbr{ʤ}ey-e&	{\hr}\tbr{ʤ}esun-e	\\	
Makatian	&	{\hr}\tbr{ʤ}im&	n-\tbr{ʤ}ip	&	\tbr{ʤ}iny-e		&	{\hr}\tbr{ʤ}u-n		&	\tbr{ʤ}uʤy-e			&	\tbr{ʤ}ew-e			&	{\hr}\tbr{ʤ}eye	&	{\hr}\tbr{ʤ}esun-e	\\	
		\lspbottomrule
	\end{tabular}
		\begin{tablenotes}\tnsize
			\item[†]{Intermediate forms from PMP *qalu-hipan are **qalipan > **lipan > **lian.}
			\item[‡]{For **luli compare Termanu \it{luli-k}, Tetun \it{luli(-k)},
							and Kisar \it{luli} `taboo, sacred' among others \citep[238]{Edwards2021}.}
			\item[Δ]{Selaru \it{sew} is just `sun'. It has \it{hey} for `day'.}
			\item[‖]{Selaru \it{nesw} ‘mortar’ is from **ŋəsun
							(probably connected with *ləsuŋ)
							which is reflected in a number of regional languages: e.g.
							Dhao \it{ŋəʧu,} Termanu \it{nesu-k,}
							Tetun \it{nesun}, and Welaun \it{kosun}.}
		\end{tablenotes}
	\end{threeparttable}
\end{adjustbox}
\end{table}

\largerpage
After high vowels,
as well as after an initial consonant and before a high vowel
(/C{\gap}V\tsc{[+high]}),
Makatian has *l > \it{ʤ}, Seluwasan has *l > \it{ʤ, d}
and Selaru has *l = \it{l}.\footnote{
		Although the data is sparse it appears that Seluwasan
		has *l > \it{d} /V\tsc{[+high]}{\gap}*a, and *l > \it{ʤ}
		/V\tsc{[+high]}{\gap}{\#}, as well as /C{\gap}V\tsc{[+high]}.
		An additional example /V\tsc{[+high]}{\gap}{\#} is Selaru \it{-ul}, Seluwasan \it{n-ud}
		Makatian \it{n-uʤ} `cut down'.}
Examples are give in \trf{04-69}.

Before final-high vowels which reduced to a glide,
Makatian and Seluwasan have the same reflexes as before high vowels,
*l > \it{ʤ} and *l > \it{d} respectively, while Selaru has a split of *l > \it{s, l}.
Makatian subsequently lost final glides after most consonants.
(It preserves glides after voiceless plosives, e.g.
pitu > \it{itw} `seven' and *matay > \it{maty} `die'.)
Examples are given in \trf{04-70}.

\largerpage
\begin{table}\tcp{\tsc{Southwest Tanimbar} *l > \it{s -- d -- ʤ} /V{\gap}\it{u, i}}{04-68}
\begin{adjustbox}{width=\textwidth}
	\begin{threeparttable}\st{2pt}
	\begin{tabular}{llllllll}\lsptoprule
local gl.	&	skin						&	bride 						&	laugh										&	egg									&	body 								&	head								&	betel 	\\	
					&									&	price							&													&											&	hair								&											&	pepper	\\	
PMP				&	*ku\tbr{l}it		&	*bə\tbr{l}i				&	**ma\tbr{l}ip						&	*qatə\tbr{l}uR			&	*bu\tbr{l}u					&	*qu\tbr{l}u					&	**ma\tbr{l}us\su{‡}	\\	\midrule
Selaru		&	\tcg{(iblu)}		&	{\hr}he\tbr{s}y		&	\hp{n}-mahi\tbr{s}\su{†}&	\hp{*}kte\tbr{s}u		&	{\hr}hu\tbr{s}y			&	{\hr}u\tbr{s}u-khatu&	{\hr}\cg{(sar)}	\\	
Seluwasan	&	{\hr}wi\tbr{d}it&	{\hr}fe\tbr{d}in	&	\hp{n-}ma\tbr{d}up			&	\hp{*k}te\tbr{d}u-n	&	{\hr}futfu\tbr{d}u-n&	{\hr}u\tbr{d}w-e		&	{\hr}ma\tbr{d}us	\\	
Makatian	&	{\hr}u\tbr{ʤ}it	&	{\hr}he\tbr{ʤ}in-e&	n-ma\tbr{ʤ}up						&	\hp{*k}te\tbr{ʤ}u-n	&	{\hr}hushu\tbr{ʤ}u-n&	{\hr}u\tbr{ʤ}u-ne		&	{\hr}ma\tbr{ʤ}us-e	\\	
		\lspbottomrule
	\end{tabular}
		\begin{tablenotes}\tnsize
			\item[†]{Selaru \it{-mahis} ‘laugh’ has consonant metathesis from earlier **masih.}
			\item[‡]{For **malus compare Tetun \it{malus},
							Meto \it{manus,} Kisar \it{maluh}
							‘betel pepper’ \citep[243f]{Edwards2021}.}
		\end{tablenotes}
	\end{threeparttable}
\end{adjustbox}
\end{table}

\begin{table}\tcp{\tsc{Southwest Tanimbar} *l > \it{ʤ -- ʤ, d -- s} /V\tsc{[+high]}{\gap}, /C{\gap}}{04-69}
	\begin{threeparttable}\st{3.9pt}
	\begin{tabular}{lllll|ll}\lsptoprule
local gl.	&	moon	&	vagina	&	maggot	&	gold	&	calm (sea)	&	ear	\\	
PMP	&	*bu\tbr{l}an	&	*ti\tbr{l}a	&	*qu\tbr{l}əj	&	*bu\tbr{l}aw-an	&	*ma-\tbr{l}inaw	&	*tə\tbr{l}iŋa	\\	\midrule
Selaru	&	{\hr}hu\tbr{l}	&	\cg{(abu)}	&	{\hr}u\tbr{l}	&	{\hr}b\tbr{l}yawan	&	\hp{*na-}m\tbr{l}in	&	{\hr}t\tbr{l}ia	\\	
Seluwasan	&	{\hr}fwi\tbr{ʤ}an	&	{\hr}ti\tbr{ʤ}y-e	&	{\hr}wi\tbr{d} miak	&	{\hr}b\tbr{ʤ}awan\su{†}	&	{\hr}na-m\tbr{d}in	&	{\hr}k\tbr{d}in-e	\\	
Makatian	&	{\hr}hu\tbr{ʤ}an	&	{\hr}ti\tbr{ʤ}a-n	&	{\hr}u\tbr{ʤ} miak	&	\cg{(mas-e)}	&	{\hr}na-m\tbr{ʤ}in	&	{\hr}k\tbr{ʤ}ina-n	\\	
		\lspbottomrule
	\end{tabular}
		\begin{tablenotes}\tnsize
			\item[†]{Seluwasan *bulaw-an > \it{bʤawan} `gold'
							appears to have *l > \it{ʤ} conditioned
							by the previous high vowel
							which was subsequently lost.
							*b = \it{b} is irregular in both
							Seluwasan and Selaru.}
		\end{tablenotes}
	\end{threeparttable}
\end{table}

\begin{table}\tcp{\tsc{Southwest Tanimbar} *l > \it{s, l -- d -- ʤ} /{\gap}G}{04-70}
\begin{adjustbox}{width=\textwidth}
	\begin{threeparttable}\st{2pt}
	\begin{tabular}{lllllllll}\lsptoprule
local gl.	&	torch							&	pestle				&	wind								&	taboo						&	rope							&	eight						&	three					&	choose	\\	
PMP				&	*su\tbr{l}uq			&	*qahə\tbr{l}u	&	**ŋə\tbr{l}u\su{†}	&	**lu\tbr{l}i		&	*ta\tbr{l}ih			&	*wa\tbr{l}u			&	*tə\tbr{l}u		&	*pi\tbr{l}iq	\\	\midrule
Selaru		&	{\hr}su\tbr{s}w		&	a\tbr{s}w			&	\hp{n}e\tbr{s}w			&	(h)asu\tbr{s}y	&	{\hr}ta\tbr{s}y		&	{\hr}wa\tbb{l}w	&	enate\tbb{l}w	&	\hp{n}-i\tbb{l}y	\\	
Seluwasan	&	{\hr}su\tbr{d}w-e	&	a\tbr{d}w-e		&	ne\tbr{d}w-e				&	du\tbr{d}y			&	{\hr}ta\tbr{d}y-e\su{‡}	&	{\hr}wa\tbr{d}w	&	rte\tbr{d}w		&	\hp{n}-i\tbr{d}y	\\	
Makatian	&	{\hr}su\tbr{ʤ}w-e	&	a\tbr{ʤ}w-e		&	ne\tbr{ʤ}w-e				&	ʤu\tbr{ʤ}y-e		&	{\hr}ta\tbr{ʤ}y-e	&	{\hr}wa\tbr{ʤ}	&	rte\tbr{ʤ}		&	n-i\tbr{ʤ}	\\	
		\lspbottomrule
	\end{tabular}
		\begin{tablenotes}\tnsize
			\item[†]{Compare Dhao \it{ŋəlu} ‘wind’ among many others (see \srf{15-04-02-02}).}
			\item[‡]{/ʤ/ in Seluwasan \it{tady-e} `rope' has an alternate
							form \it{taʤ-e} which appears to be
							due to a synchronic /dy/ {\ra} /ʤ/ /{\gap}V.
							Compare \it{dudy} `taboo, sacred' with \it{ton duʤ-e} `sacred place'.}
		\end{tablenotes}
	\end{threeparttable}
\end{adjustbox}
\end{table}

In other environments, that is before \it{a} and \it{o}
and word finally after \it{a} (/{\gap}\it{a,o} and /\it{a}{\gap}{\#}),
all three languages have *l = \it{l}.
Examples are given in \trf{04-71}.

\begin{table}\tcp{\tsc{Southwest Tanimbar} *l = \it{l} /{\gap}\it{a,o}, /\it{a}{\gap}{\#}}{04-71}
	\begin{threeparttable}\st{6pt}
	{\st{2.6pt}\begin{tabular}{llllllll}\lsptoprule
local gl.	&	husband						&	sky									&	\hp{na-}run					&	sail							&	forest								&	hungry							&	inside	\\	
PMP				&	*\tbr{l}aki				&	*\tbr{l}aŋit				&	\hp{na}*\tbr{l}aRiw	&	*\tbr{l}ayaR			&	*ha\tbr{l}as					&	*(ma-)\tbr{l}apaR		&	*da\tbr{l}əm	\\	\midrule
Selaru		&	{\hr}\tbr{l}ai		&	{\hr}\tbr{l}ait			&	\hp{na}-\tbr{l}a		&	{\hr}\tbr{l}ar		&	{\hh}a\tbr{l}as\su{†}	&	\hp{||na-m}\tbr{l}ar&	{\hr}kra\tbr{l}a	\\	
Seluwasan	&	{\hr}\tbr{l}ai-n	&	{\hr}\tbr{l}anit-e	&	\hp{na-}\tbr{l}aw		&	{\hr}\tbr{l}aw-e	&	{\hh}a\tbr{l}as				&	\hp{||}na-m\tbr{l}ay&	{\hr}\tbr{l}arum\su{‡}	\\	
Makatian	&	{\hr}\tbr{l}ai-n	&	{\hr}\tbr{l}anit		&	na-\tbr{l}aw				&	{\hr}\tbr{l}aw-e	&	{\hh}a\tbr{l}as-e			&	\hp{*-}ta-m\tbr{l}ay&	{\hr}rya\tbr{l}am	\\	
	\midrule\midrule
	\end{tabular}}
	{\st{4.25pt}\begin{tabular}{llllllll}
local gl.	&	mouse							&	fly (n.)								&	road							&	get							&	swallow							&	sea								&	throw up\su{Δ}		\\	
PMP				&	*ba\tbr{l}abaw		&	*\tbr{l}a\tbr{l}əj			&	*za\tbr{l}an			&	*a\tbr{l}ap			&	*tə\tbr{l}ən				&	*\tbr{l}ahud			&	*\tbr{l}uaq	\\\midrule
Selaru		&	{\hr}k\tbr{l}af		&	{\hr}\tbr{l}a\tbr{l}		&	{\hr}sa\tbr{l}		&	\hp{n}-a\tbr{l}	&	\hp{n}-te\tbr{l}		&	\cg{(tasi)}				&		\\	
Seluwasan	&	{\hr}f\tbr{l}af-e	&	{\hr}\tbr{l}a\tbr{l}-e	&	{\hr}sya\tbr{l}an	&	\hp{n-}a\tbr{l}	&	\hp{n-}te\tbr{l}an	&	{\hr}\tbr{l}ot		&	n-\tbr{l}on	\\	
Makatian	&	{\hr}h\tbr{l}ah-e	&	{\hr}\tbr{l}a\tbr{l}-e	&	{\hr}sya\tbr{l}an	&	n-a\tbr{l}			&	n-te\tbr{l}an				&	{\hr}\tbr{l}ot-e	&	n-\tbr{l}ow	\\	
		\lspbottomrule
	\end{tabular}}
		\begin{tablenotes}\tnsize
			\item[†]{Selaru \it{alas} is ‘dense forest’.}
			\item[‡]{Seluwasan \it{larum} has consonant metathesis.}
			\item[Δ]{Specifically refers to a baby throwing up.}
			%\item[‖]
			%\item[¶]
		\end{tablenotes}
	\end{threeparttable}
\end{table}

Given these reflexes, it is almost certain that Selaru \it{s}
is from earlier **ʤ. *ʤ > \it{s} is well attested sound change
(also seen in \tsc{Southwest Tanimbar} *z [ʤ] > \it{s}).
Additional evidence for this comes from a number of words
in Makatian which show *ʤ > \it{s}.
These include *bulu > \it{hu\tbr{s}huʤu-n} `body hair',
**ŋəlu > \it{neʤw-e} wind, but \it{ne\tbr{s} timur} `east wind',
as well as **malip > \it{n-maʤup} `laughs' but \it{na-ma\tbr{s}pu ni} `laugh at him/her'.
We thus that propose Selaru has *l > **ʤ > \it{s}.

Similarly, we propose that Seluwasan has **ʤ > \it{d, ʤ}
rather than **d > \it{d, ʤ}.
The former solution is simpler as it means we only need to propose
three changes: (i) *l > Proto-Southwest Tanimbar **ʤ
(ii) Seluwasan **ʤ > \it{d}
and (iii) Selaru **ʤ > \it{s} (each in the appropriate environments).
However, if we were to propose *l > Proto-Southwest Tanimbar **d,
this would entail an additional change for Makatian and probably Selaru.
We thus propose the sound changes given in \qf{04-29}.
(H = high vowel and G = glide.)

{\wg\st{3pt}\tblr{>{\hspace{-\tabcolsep}}l<{\hspace{-\tabcolsep}}>{\hspace{-\tabcolsep}\enspace}l
									>{\hspace{-\tabcolsep}\space\space}l
									>{\hspace{-\tabcolsep}\space\space}l
									>{\hspace{-\tabcolsep}\space\space}lll}
\tbx{04-29}	&	\mc{6}{l}{Changes affecting *l in \tsc{Southwest Tanimbar} (Scenario 1)}	\\
						&	a.	&	{\hr}*l	&	>	&	**ʤ						&	/{\gap}H, e							&	Proto-Southwest Tanimbar	\\
						&	b.	&	{\hr}*l	&	>	&	**ʤ, **l			&	/{\gap}G{\#}						&	Proto-Southwest Tanimbar	\\
						&	c.	&	**l			&	>	&	**ʤ 					&	/{\gap}G{\#}, /H{\gap}	&	Proto-Makatian-Seluwasan	\\
						&	d.	&	**ʤ			&	>	&	\hp{**}\it{d}	&	/{\gap}H, /{\gap}(G){\#}&	Seluwasan	\\
						&	e.	&	**ʤ			&	>	&	\hp{**}\it{s}	&													&	Selaru	\\
\end{tabular}\mbox{}\\}

It is possible that the split at (\ref{ex:04-29}b)
is due to Selaru having already
shifted some (but not all) final high vowels to glides.
Similarly, the conditioning environment /{\gap}G{\#}
for Makatian and Seluwasan (in \ref{ex:04-29}c and \ref{ex:04-29}d)
can be eliminated by proposing that the shift of final high
vowels to glides took place after the changes affecting *l.
This scenario is outlined in \qf{04-29b}.
Under this scenario the change of final high vowels
to glides was completed after the break up of the proto-language.

{\wg\st{3pt}\tblr{>{\hspace{-\tabcolsep}}l<{\hspace{-\tabcolsep}}>{\hspace{-\tabcolsep}\enspace}l
									>{\hspace{-\tabcolsep}\space\space}l
									>{\hspace{-\tabcolsep}\space\space}l
									>{\hspace{-\tabcolsep}\space\space}lll}
\tbx{04-29b}	&	\mc{6}{l}{Changes affecting *l in \tsc{Southwest Tanimbar} (Scenario 2)}	\\
							&	a.	&	{\hr}*H	&	>	&	\hp{**}G, *H	&	/{\gap}{\#}	&	Selaru	\\
							&	b.	&	{\hr}*l	&	>	&	**ʤ						&	/{\gap}H, e	&	\tsc{Southwest Tanimbar}	\\
							&	c.	&	{\hr}*l	&	>	&	**ʤ 					&	/H{\gap}		&	\tsc{Makatian-Seluwasan}	\\
							& d.	&	**ʤ			&	>	& \hp{**}\it{d}	&	/{\gap}H		& Seluwasan \\
							& e.	&	**ʤ			&	>	& \hp{**}\it{s}	&							& Selaru \\
							&	f.	&	{\hr}*H	&	>	&	\hp{**}G			&	/{\gap}{\#}	&	\tsc{Southwest Tanimbar}	\\
\end{tabular}\mbox{}\\}

\largerpage
Supporting evidence for \tsc{Southwest Tanimbar} comes from shared
*j > {\0}, *b > **f (with subsequent **f > \it{h} in Makatian
and Selaru),\footnote{
		There is also one instance of word-final
		*b in our data: *ma-huab > Makatian \it{n-ap},
		Seluwasan \it{n-mwap}, and Selaru \it{-maf} ‘yawn’.}
*ə > \it{e} in non-final syllables,
and *d > \it{r} (\trf{04-65}).
Examples of *j > {\0} are given in \trf{04-72}.
%and examples of *b > **f are given in \trf{04-73}.
Note that Makatian and Seluwasan reflexes of *ŋajan > \it{naw-n} ‘name’
have *j > \it{w}, parallel to changes affecting *R (discussed below).

\begin{table}\tcp{\tsc{Southwest Tanimbar} *j}{04-72}
	\begin{threeparttable}\st{5pt}
	\begin{tabular}{lllllll}\lsptoprule
local gl.	&	how m.?			&	sibling\su{†}	& gall(bladder)	& pitfall			&	sun, day 				&	name							\\
PMP				&	*pi\tbr{j}a	&	*hua\tbr{j}i	&	*qapə\tbr{j}u	& *su\tbr{j}a	&	*qalə\tbr{j}aw	&	*na\tbr{j}an			\\ \midrule
Selaru		&	\cg{(enai)}	&	{\hr}wai-n		&	ew						&							&	{\hr}sew\su{Δ}	&	\hp{*n}any				\\
Seluwasan	&	{\hr}ii			&	{\hr}wai-n		& ew-n					&	suw-e				&	{\hr}ʤew-e			&	{\hr}na\tbr{w}-e	\\
Makatian	&	{\hr}ii			&	{\hr}wai-n		&	ew-ne					&	suw-e\su{‡}	&	{\hr}ʤew-e			&	{\hr}na\tbr{w}-n	\\
		\lspbottomrule
	\end{tabular}
		\begin{tablenotes}\tnsize
			\item[†]{Selaru \it{wai-n} is `sibling, friend;
							%used by same sex siblings; to same sex friends
							%(when relation is not clear) or
							%any same sex person of the same village or community
							%(when referring to them when talking with someone else);
							%also used to refer to anyone of the opposite sex who
							potential marriage partner'.}
							%also used after marriage to refer to all
							%same sex relations of the partner
							%(all wife's sisters; all husband's brothers).'}
			\item[‡]{Medial \it{w} in Makatian and Seluwasan \it{suw-e} `pitfall trap'
							could be due to *j > \it{w}, or it could be an epenthetic
							transitional glide between two vowels.}
			\item[Δ]{\it{w} in reflexes of *qaləjaw is
							probably a reflex of *-aw or *w,
							%(perhaps *-aw > \it{w} /V{\gap} after *j > {\0}),
							and not *j/*R > \it{w} as Selaru does not have *R > \it{w}.
							However, *-aw > {\0} /C{\gap} in all three languages.
							%See reflexes of *ma-linaw `calm (sea)' (\trf{04-69}, \prf{tab:04-69})
							%*balabaw `mouse, rat' (\trf{04-71}, \prf{tab:04-71}),
							%and *bətaw (\trf{04-72}, \prf{tab:04-72}).
							Selaru \it{sew} is just `sun'.
							}
		\end{tablenotes}
	\end{threeparttable}
\end{table}

%\begin{table}\tcp{\tsc{Southwest Tanimbar} *b > **f}{04-73}
%\begin{adjustbox}{width=\textwidth}
	%\begin{threeparttable}\st{2pt}
	%\begin{tabular}{llllllll}\lsptoprule
%local gl.	&	woman	&	stone	&	lips	&	head hair	&	mouth	&	new	&	mouse\\
%PMP	&	*\tbr{b}ətaw	&	*\tbr{b}atu	&	*\tbr{b}i\tbr{b}iR	&	*\tbr{b}uhək	&	*\tbr{b}a\tbr{b}aq	&	*\tbr{b}aqəRu	&	*\tbr{b}ala\tbr{b}aw\\ \midrule
%Selaru	&	{\hr}wam\tbr{f}wet	&	{\hr}\tbr{h}atw	&	{\hr}\tbr{h}i\tbr{h}i	&	{\hr}\tbr{h}uka-	&	{\hr}\tbr{h}a\tbr{h}a-\su{†}	&	{\hr}\tbr{h}ar{\tl}harw	&	{\hr}kla\tbr{f}\\
%Seluwasan	&	{\hr}\tbr{f}et-e	&	{\hr}\tbr{f}atw-e	&	{\hr}\tbr{f}i\tbr{f}i-ne	&	{\hr}\tbr{f}uk-e	&	{\hr}\tbr{f}a\tbr{f}a-n	&	{\hr}\tbr{f}aw	&	{\hr}\tbr{f}la\tbr{f}-e\\
%Makatian	&	{\hr}\tbr{h}et-e	&	{\hr}\tbr{h}atw-e	&	{\hr}\tbr{h}i\tbr{h}i-ne	&	{\hr}\tbr{h}uk-e	&	{\hr}\tbr{h}a\tbr{h}a-n	&	{\hr}\tbr{h}aw	&	{\hr}\tbr{h}la\tbr{h}-e\\
		%\lspbottomrule
	%\end{tabular}
		%\begin{tablenotes}\tnsize
			%\item[†]	{Selaru \it{haha-} is ‘inside of mouth’.}
		%\end{tablenotes}
	%\end{threeparttable}
%\end{adjustbox}
%\end{table}

Within \tsc{Southwest Tanimbar}, Makatian
and Seluwasan belong together.
One sound change shared between them which sets them apart from
Selaru are the changes affecting *R.
Makatian and Seluwasan share a split of *R > {\0}, \it{r, w, y}
while Selaru usually has initial/medial *R > \it{r},
and word-final *R > {\0}, \it{r}.
Examples of *R > \it{r} in all three languages are given in \trf{04-74}.

\begin{table}\tcp{\tsc{Southwest Tanimbar} *R > \it{r}}{04-74}
\begin{adjustbox}{width=\textwidth}
	\begin{threeparttable}\st{2pt}
	\begin{tabular}{llllllll}\lsptoprule
local gl.	&	west	&	canarium	&	octopus	&	conch	&	east	&	spit	&	dove	\\
PMP				&	*haba\tbr{R}at			&	*kana\tbr{R}i\su{†}	&	*ku\tbr{R}ita			&	*tabu\tbr{R}i(q)		&	*timu\tbr{R}			&	*bu\tbr{R}ah			&	**lakatə\tbr{R}u\su{‡}	\\ \midrule
Selaru		&	\hp{*ha}ha\tbr{r}at	&	{\hr}kna\tbr{r}y		&	\cg{(amana)}			&	{\hr}htu\tbr{r}y		&	{\hr}timu\tbr{r}	&	\cg{(-ohuk)}			&	{\hr\hr}manu lakte\tbr{r}w	\\
Seluwasan	&	\hp{*ha}fa\tbr{r}at	&	{\hr}kna\tbr{r}y-e	&	{\hr}k\tbr{r}ite	&	{\hr}sbu\tbr{r}y-e	&	{\hr}timu\tbr{r}	&	{\hr}n-fu\tbr{r}a	&	{\hr\hr}tuk\tbr{r}u-e	\\
Makatian	&	\hp{*ha}ha\tbr{r}at	&	{\hr}nya\tbr{r}y-e	&	{\hr}k\tbr{r}ite	&	{\hr}sbu\tbr{r}y-e	&	{\hr}timu\tbr{r}	&	{\hr}hwi\tbr{r}a	&	{\hr\hr}lkot\tbr{r}u-e	\\
		\lspbottomrule
	\end{tabular}
		\begin{tablenotes}\tnsize
			\item[†]{\citet{BlustTrussel2020} reconstruct PMP *kanari with medial *r.
							However, the words cited in favour of *r rather than *R are
							suspect loans from Malay \it{kənari} and/or
							the reflexes of *r as opposed to *R
							are not firmly established. \citet[315f]{Ross2008b}
							reconstructs PCEMP **ka(nŋ)aRi and \tsc{POc} *(ka)ŋaRi.}
			\item[‡]{For **lakatəRu compare  Tetun \it{lakateu}, Waima{\Q}a \it{lakateu},
							 Teun \it{lakateri}, and Leti \it{laktyeru} `turtle dove'.}
			%\item[Δ]
			%\item[‖]
			%\item[¶]
		\end{tablenotes}
	\end{threeparttable}
\end{adjustbox}
\end{table}

Examples of *R > {\0} in Makatian and Seluwasan where
Selaru has *R > \it{r} are given in \trf{04-75}.
There is also one case of *r in \trf{04-75}.
Examples in which all languages
have *R > {\0} /{\gap}{\#} are given in \trf{04-76}.
Many words with final *R > {\0} are
nouns which usually take a genitive suffix. 

\begin{table}\tcp{\tsc{Southwest Tanimbar} *R > \it{r} -- {\0} -- {\0}}{04-75}
%\begin{adjustbox}{width=\textwidth}
	\begin{threeparttable}\st{2pt}
	\begin{tabular}{lllllll}\lsptoprule
local gl.	&	stingray	&	new	&	bone	&	candlenut	&	ladder	&	calcium	\\	
PMP	&	*pa\tbr{R}ih	&	*baqə\tbr{R}u		&	*du\tbr{R}i	&	*kami\tbr{r}i	&	*ha\tbr{R}əzan\su{†}	&	*qapu\tbr{R}	\\	\midrule	
Selaru	&	\hp{*p}a\tbr{r}y	&	{\hr}ha\tbr{r}ha\tbr{r}w	&	{\hr}lu\tbr{r}y	&	{\hr}kmu\tbr{r}y	&	{\hh}a\tbr{r}es	&	{\hq}au\tbr{r}	\\	
Seluwasan	&	\hp{*p}ay-e	&	{\hr}faw		&	{\hr}ruy-e	&	\hp{*k}muy-e	&	{\hh}esan	&	{\hq}aw-e	\\	
Makatian	&	\hp{*p}ay-e	&	{\hr}haw-e		&	{\hr}ruye	&	\hp{*k}muye	&	{\hh}esan	&	{\hq}aw-e	\\
		\lspbottomrule
	\end{tabular}
		\begin{tablenotes}\tnsize
			\item[†]{Makatian and Seluwasan preserve the second
							two syllables of *haRəzan, while Selaru preserves the first two.}
			%%\item[‡]
			%\item[Δ]
			%\item[‖]
			%\item[¶]
		\end{tablenotes}
	\end{threeparttable}
%\end{adjustbox}
\end{table}

\begin{table}\tcp{\tsc{Southwest Tanimbar} *R > {\0}}{04-76}
\begin{adjustbox}{width=\textwidth}
	%\begin{threeparttable}
	\st{2pt}\begin{tabular}{lllllllll}\lsptoprule
local gl.	&	hundred	&	man	&	bow (n.)	&	penis	&	egg	&	tail	&	lips	&	root	\\	
PMP	&	*\tbr{R}atus	&	*ma\tbr{R}uqanay	&	*busu\tbr{R}	&	*lasə\tbr{R}	&	*qatəlu\tbr{R}	&	*iku\tbr{R}	&	*bibi\tbr{R}	&	*waka\tbr{R}	\\	\midrule	
Selaru	&	{\hr}atw	&	{\hr}wamwany	&	{\hr}husw	&	{\hr}lasa	&	{\hr}ktesu	&	{\hr}iku	&	{\hr}hihi	&	\cg{(kawan)}	\\	
Seluwasan	&	{\hr}asut \it{(m.)}	&	{\hr}mwan-e	&	{\hr}fusw-e	&	{\hr}las-e	&	{\hr}tedu-n	&	{\hr}iku-n	&	{\hr}fifi-n	&	{\hr}wa-ne	\\	
Makatian	&	{\hr}asut \it{(m.)}	&	{\hr}mwane	&	{\hr}husw-e	&	{\hr}lasa-n	&	{\hr}teʤu-n	&	{\hr}iku-ne	&	{\hr}hihi-n	&	{\hr}wa-ne	\\
		\lspbottomrule
	\end{tabular}
		%\begin{tablenotes}\tnsize
			%\item[†]
			%%\item[‡]
			%%\item[Δ]
			%%\item[‖]
			%%\item[¶]
		%\end{tablenotes}
	%\end{threeparttable}
\end{adjustbox}
\end{table}

The other reflex of *R in Makatian and Seluwasan is \it{w, y}.
This usually (though not universally) corresponds to Selaru *R > \it{r}.
Several forms which appear to attest *R > \it{w} are
ambiguous between *R > {\0} with subsequent
epenthesis of a glide after a high vowel, such as
*niuR > Makatian, Seluwasan [ˈnuwɛ] (possibly /nuw-e/ or /nu-e/), Selaru \it{nur} `cononut'.
%or *uRat > Makatian, Seluwasan [ˈuwat̪] (possibly /uat-e/ or /uwat-e), Selaru \it{urta} `vein'.
Similarly, some apparent cases of *R > \it{w, y} may be due
to *R > {\0} with the final glide from a final high vowel or glide,
such as %*wahiyəR > Makatian and Seluwasan \it{wey-e}, Selaru \it{wer} `water' (possibly from intermediate **waiR) or
*laRiw > Makatian and Seluwasan \it{law}, Selaru \it{-la} `run'.
Probable unambiguous cases of *R > \it{w, y} are given in \trf{04-77}.

\begin{table}\tcp{\tsc{Southwest Tanimbar} *R > \it{w, y} -- \it{w, y} -- \it{r}, {\0}}{04-77}
	\begin{threeparttable}\st{6pt}
	\begin{tabular}{llllll}\lsptoprule
local gl.	&	\hp{n}come				&	sail (n.)					&	blood						&	hungry							&	red	\\	
PMP				&	\hp{n}*ma\tbr{R}i	&	*laya\tbr{R}			&	*da\tbr{R}aq		&	*(ma-)lapa\tbr{R}		&	*ma-i\tbr{R}aq	\\	\midrule
Selaru		&	\hp{n-}mai				&	{\hr}la\tbr{r}		&	{\hr}la\tbr{r}a	&	\hp{*na-m}la\tbr{r}	&	{\hr}mer{\tl}me\tbr{r}	\\	
Seluwasan	&	\hp{n-}ma\tbr{w}	&	{\hr}la\tbr{w}-e	&	{\hr}ra\tbr{w}-e&	{\hr}na-mla\tbr{y}				&	{\hr}me{\tl}me\tbr{y}-e	\\	
Makatian	&	n-ma\tbr{w}				&	{\hr}la\tbr{w}-e	&	{\hr}ra\tbr{w}-e&	\hp{*|}ta-mla\tbr{y}	&	{\hr}beve\tbr{y}-e\su{†}	\\	
		\lspbottomrule
	\end{tabular}
		\begin{tablenotes}\tnsize
			\item[†]{Makatian \it{bevey-e} `red' is probably
							a reduplicated form from **bebey.
							If it is from *ma-iRaq
							it has irregular *m > \it{b}.}
		\end{tablenotes}
	\end{threeparttable}
\end{table}

Another difference between Selaru
and \tsc{Makatian-Seluwasan}
is that the latter usually have *ŋ > \it{n}
while Selaru has a split of *ŋ > \it{n,} {\0}.
Examples in which Selaru has *ŋ > \it{n} are given in \trf{04-78},
and examples for which Selaru has *ŋ > {\0} are given in \trf{04-79}.

\begin{table}\tcp{\tsc{Southwest Tanimbar} *ŋ > \it{n}}{04-78}
%\begin{adjustbox}{width=\textwidth}
	\begin{threeparttable}\st{4.25pt}
	\begin{tabular}{lllllll}\lsptoprule
local gl.	&	flower	&	branch	&	cold	&	tooth	&	soul	&	rudder	\\	
PMP	&	*bu\tbr{ŋ}a	&	*sa\tbr{ŋ}a	&	*ma-di\tbr{ŋ}in	&	*\tbr{ŋ}isi	&	*suma\tbr{ŋ}əd	&	*quli\tbr{ŋ}	\\	\midrule
Selaru	&	{\hr}khu\tbr{n}a	&	{\hr}ksa\tbr{n}a	&	{\hr}m(b)ri\tbr{n}\su{†}	&	{\hr}\tbr{n}isi	&	{\hr}smwak\tbr{n}a	&	{\hr}wili\tbr{n}	\\	
Seluwasan	&	{\hr}fu\tbr{n}a-n	&	{\hr}sa\tbr{n}a-n	&	{\hr}bri\tbr{n}in	&	{\hr}\tbr{n}isy-e	&	{\hr}smwa\tbr{n}t-ane	&	{\hr}wiʤ-e	\\	
Makatian	&	{\hr}hu\tbr{n}a-ne	&	{\hr}sa\tbr{n}a-n	&	{\hr}mri\tbr{n}in	&	{\hr}\cg{(nii-n)}	&	{\hr}smwa\tbr{n}t-an	&	{\hr}wiʤ-e	\\	
		\lspbottomrule
	\end{tabular}
		\begin{tablenotes}\tnsize
			\item[†] {Selaru \it{mrin, mbrin} ‘cold’ is ambiguous between
								showing *ŋ > {\0} and retention of the final *n
								(*ma-diŋin > **madiin > \it{mrin}),
								or *ŋ > \it{n} with loss of the final vowel and consonant.}
		\end{tablenotes}
	\end{threeparttable}
%\end{adjustbox}
\end{table}

\begin{table}\tcp{\tsc{Southwest Tanimbar} *ŋ > {\0} -- \it{n} -- \it{n}}{04-79}
%\begin{adjustbox}{width=\textwidth}
	\begin{threeparttable}
	\st{4pt}\begin{tabular}{llllllll}\lsptoprule
local gl.	&	name	&	wind	&	year\su{†}	&	sky	&	swim	&	ear	&	dugong	\\	
PMP	&	*\tbr{ŋ}ajan	&	**\tbr{ŋ}əlu	&	*ha\tbr{ŋ}in	&	*la\tbr{ŋ}it	&	*na\tbr{ŋ}uy	&	*təli\tbr{ŋ}a	&	*duyu\tbr{ŋ}	\\	\midrule
Selaru	&	\hp{*n}any	&	\hp{*n}esw	&	{\hh}ain	&	{\hr}lait	&	-naw	&	{\hr}tlia	&	{\hr}rui	\\	
Seluwasan	&	{\hr}\tbr{n}aw-e	&	{\hr}\tbr{n}edw-e	&	{\hh}a\tbr{n}in	&	{\hr}la\tbr{n}it-e	&	{\hr}na\tbr{n}y	&	{\hr}kdi\tbr{n}-e	&	{\hr}rui\tbr{n}-e	\\	
Makatian	&	{\hr}\tbr{n}aw-n	&	{\hr}\tbr{n}eʤw-e	&	{\hh}a\tbr{n}in	&	{\hr}la\tbr{n}it	&	{\hr}nya\tbr{n}w	&	{\hr}kʤi\tbr{n}a-n	&	{\hr}ru\tbr{n}-e	\\	
		\lspbottomrule
	\end{tabular}
		\begin{tablenotes}\tnsize
			\item[†]{PMP *haŋin is `wind'.}
			%\item[‡]
			%\item[Δ]
			%\item[‖]
			%\item[¶]
		\end{tablenotes}
	\end{threeparttable}
%\end{adjustbox}
\end{table}
